% \chapter{Introduction}
\chapter{Superhydrophobicity-Enabled Elastic Interfacial Effect Field Transistor Based on Multiscale Interfacial Phenomena}
\label{ch:small}
\renewcommand*\imgdir{img/ch-small/}

% \dictum[Wolfgang Pauli]{%
  % God made the bulk;\\Surfaces were invented by the devil
  % }%
\worktodo{find the quote or quit}

\vspace{1em}

\chapterabstract{Part of this chapter appears in the following journal
  article: Tian, T., Sharma, C., Navanshu, A., Varga, M., Lee, Y.-T.,
  Chiu, Y.-C. \& Shih, C.-J. An Elastic Interfacial Transistor Enabled
  by Superhydrophobicity. Small 14, 1804006 (2018).  }



% \newpage{}
% \begin{abstract}
%   % Having been the central component of the semiconductor industry for
%   % more than 70 years, the field-effect transistors (FETs) still
%   % attract substantial interest and development. One long-existing
%   % challenge is to incorporate mechanical responsiveness and durability
%   % into FETs, which requires tedious material design and may compromise
%   % device performance. Here we introduce an alternative approach, the
%   % interfacial field-effect transistor (IFET), combining multiscale
%   % interfacial physical phenomena, to address such problem. The
%   % mechanical responsiveness of IFET is harnessed from a conductive
%   % Cassie-Baxter sessile droplet on superhydrophobic semiconducting
%   % nanowires, formed by templated growth on the interface of a
%   % two-dimensional electron gas (2DEG). The extremely low elastic
%   % modulus of the droplet, much lower
%   % than any existing rubbers, ensures an excellent elastic sensiting
%   % limit down to \textless{}10 Pa. The partial penetration of electric field
%   % through the 2DEG modulates the carrier profile at the NW/2DEG
%   % interface and a thermal-tunable current on/off ratio exceeding
%   % $3\times{}10^{4}$. This study demonstrates a versatile platform that
%   % bridges multiple macroscopic interfacial phenomena with
%   % nanoelectronic responses, which can be used to facilely extend the
%   % functionality of semiconductor components.
%   Enabling mechanical responsiveness in field-effect transistors
%   (FETs) offers new technological opportunity beyond the reach of
%   existing platforms. Here we propose a new force-sensing concept by
%   controlling the wettability of a semiconductor surface, referring to
%   the interfacial field-effect transistors (IFETs). We design an IFET
%   made by superhydrophobic semiconductor nanowires (NWs) sandwiched
%   between a layer of two-dimensional electron gas (2DEG) and a
%   conductive Cassie-Baxter (CB) sessile droplet. Following the
%   hydrostatic deformation of the CB droplet upon mechanical stress, we
%   report an extremely small elastic modulus of 820 pascals vertical to
%   the substrate plane, or $\sim{}$100 times softer than Ecoflex rubbers,
%   enabling an excellent stress detection limit down to \textless{}10
%   pascals and a stress sensitivity of 36 kPa$^{-1}$. The IFET exhibits
%   an on/off current ratio exceeding $3\times{}10^{4}$, as the carrier
%   density profile at the NW/2DEG interface is modulated by a
%   partially-penetrated electrostatic field. This study demonstrates a
%   versatile platform that bridges multiple macroscopic interfacial
%   phenomena with nanoelectronic responses.
% \end{abstract}

% % \begin{tocentry}
% %   \vfill
% %   \includegraphics[width=0.98\linewidth]{raw_img/TOC.pdf}
% %   \vfill
% % \end{tocentry}

% \linenumbers

Recent advances in semiconductor materials and physics have
driven the development of field-effect transistors (FETs) for a
variety of applications \cite{Torsi_2013_rev,Ben_Sasson_2014}.  Enabling
mechanical durability and responsiveness in FETs offers new
technological opportunities beyond the reach of existing platforms
\cite{Someya_2004,Oh_2016,Shin_2017}. It is desirable to intrinsically
incorporate the responsiveness in semiconductor by rational materials
design \cite{Oh_2016} but often compromises the device performance and
processability \cite{Lee_2018}. For example, the elasticity, the ability
of a soft material to return to its original shape when the deforming
stress is released, requires an entropic force acting within the
material, fundamentally impeding carrier transport
\cite{OConnor_2011_strain_P3HT}. As a result, so far there are only few
reports that have successfully demonstrated stress-sensing FETs using
intrinsically flexible semiconductors \cite{Oh_2016}. On the other hand,
incorporating transistor components into a mechanically flexible
matrix also offers opportunities for stress-sensing FETs. A widely
used strategy involves conductive elastic polymer as the
pressure-sensitive unit
\cite{Someya_2004,Sekitani_2009,Kaltenbrunner_2013,Takei_2010}, of which
the sensitivity is limited by the properties of the polymer-based
components. Microstructures embedded in an elastomer enhances the
mechanical durability \cite{Jang_2015}, and can be used to improve the
flexibility and responsiveness of flexible rubber-based FETs
\cite{Mannsfeld_2010,Schwartz_2013}. Recently, highly sensitive FETs
using air as the compressible dielectric layer has as also been
proposed \cite{Zang_2015,Shin_2017}. The elasticity of the polymer-based
FETs essentially depends on the elastic modulus of the flexible
material, which is typically in the range of \(10^{1}\sim{}10^{3}\) kPa
\cite{Amjadi_2016}. Here we present a conceptually
different approach, referring to the interfacial field-effect
transistor (IFET), in which the mechanical responsiveness is
controlled by the wettability of a semiconductor surface.

Figure \ref{fig:main-1} presents the device architecture of the proposed
IFET. We consider a quantum capacitor device \cite{Luryi_1988}, in which
a semiconductor layer is sandwiched between a conductive drain (D) and
a 2DEG source (S) electrodes. The ultrathin and
low-quantum-capacitance (low \(C_{\mathrm{Q}}\)) nature of 2DEG allows
partial penetration of an electrostatic field exerted by the
underlying gate (G) electrode \cite{Shih2015Partially,Tian_2016},
modulating the carrier density in semiconductor that controls
drain-to-source current (\(I_{\mathrm{DS}}\)) (Figure
\ref{fig:main-1}a).  In addition, the fact that \(I_{\mathrm{DS}}\)
passes through two semiconductor-electrode interfaces maximizes the
response to the interfacial properties.

Next, we sought to grow semiconductor NWs directly on the 2DEG,
forming a superhydrophobic surface at the Cassie-Baxter (CB) state
\cite{Cassie_1944}. A conductive sessile droplet is then placed on top
of the NWs, acting as the D electrode, with the radius of contact,
\(R\). Accordingly, applying a compressive stress \(\sigma\) on top
of the sessile leads to a horizontal displacement \(\Delta x\) that
advances the contact line, with an apparent advancing contact angle
\(\theta_{\mathrm{adv}}^{*}\) (Figure \ref{fig:main-1}b left). When the
stress is released, the contact line moves back towards its
equilibrium shape, with an apparent receding contact angel of
\(\theta_{\mathrm{rec}}^{*}\) (Figure \ref{fig:main-1}b right), enabling the
reversibility. We hypothesized that the electrical response during the
process results from: (i) an increase of the droplet contact area
\(\Delta A \approx 2 \pi R \Delta x\) at the macroscopic level and (ii)
an increase of penetration length \(\Delta l \approx \sigma
r^{2}/2\gamma_{\mathrm{L}}\) that effectively reduces the semiconductor
channel length in the NWs (see Figure  \ref{fig:main-1}b bottom), where
\(\gamma_{\mathrm{L}}\) is the surface tension of liquid and \(r\) is the
average radius between NWs.  In order to realize the concept,
monolayer graphene (Gr), a 2DEG material with the smallest quantum
capacitance near the Dirac point \cite{Tian_2016}, is used. As for the
semiconductor material, it is well-known that the covalently bonded
semiconductors generally possess a high surface energy, owing to the
high surface polarity \cite{Azimi_2013_wetting_RO} and dangling bonds
\cite{Zhang_2004_dangling}, presumably leading to the Wenzel state. To
this end, the organic semiconductor solids containing the
\(\pi\)-conjugated molecules, with the molecular polarizability
considerably reduced by perfluoration, were considered. We have
managed to grow NWs of Copper(II) hexadecafluorophthalocyanine
(F\(_{\text{16}}\)CuPc), a n-type organic semiconductor with high air stability
\cite{Bao_1998}, on a sheet of monolayer graphene transferred onto a
SiO\(_{\text{2}}\)/Si substrate (details see Section S1 and 
Figure S2 in the Supporting Information), as shown in Figure \ref{fig:main-2}a.

The F\(_{\text{16}}\)CuPc NWs have high aspect ratios, with diameters and heights
of 30 - 60 and several hundred nanometers, respectively. We notice that
the NW morphology is similar to the F\(_{\text{16}}\)CuPc films templated by
perylene-3,4,9,10-tetracarboxylic-3,4,9,10-dianhydride (PTCDA)
\cite{Yang_2009_F16_PTCDA} or gold nanoparticles
\cite{Mbenkum_2006_F16_1D}, presumably sharing the same growth
mechanism. Indeed, the strong \(\pi-\pi\) interactions between graphene
and F\(_{\text{16}}\)CuPc basal planes lower the critical free energy of
heterogeneous nucleation; in combination with the weak van der Waals
(vdW) interactions between the perfluorinated molecular edges that
limit the lateral growth, it explains why the NW structure is
thermodynamically favored. The structural driving force resulting from
the subtle asymmetry between in-plane and out-of-plane interactions,
however, disappears when the substrate-molecule interactions are not
sufficiently strong \cite{Tian_2017}, as observed in the thin-film
structure grown on bare SiO\(_{\text{2}}\) (Figure \ref{fig:main-2}a). More analysis
about the morphologies of F\(_{\text{16}}\)CuPc NW and F\(_{\text{16}}\)CuPc/SiO\(_{\text{2}}\) samples
see Section S2 and Figures S3-S9 in the Supporting Information.

Figure \ref{fig:main-2}b shows the grazing incidence X-ray diffraction
(GIXD) patterns for the F\(_{\text{16}}\)CuPc NW and F\(_{\text{16}}\)CuPc/SiO\(_{\text{2}}\)
samples. The F\(_{\text{16}}\)CuPc/SiO\(_{\text{2}}\) sample has an intense peak along the
\(q_{\mathrm{z}}\) axis, labeled as \(\alpha\) (\emph{d}-spacing = 1.470
nm). It corresponds to the (002) plane in the F\(_{\text{16}}\)CuPc single
crystal \cite{Yang_2009_F16_PTCDA} that is perpendicular to the
substrate normal vector, suggesting a dominant ``edge-on'' orientation
(Figure \ref{fig:main-2}b bottom left). It is also endorsed by the
position of the (\(\bar{1}\)22) plane peak near the \(q_{\mathrm{xy}}\)
axis (\(\beta\)) \cite{Yoon_2010_crystal_F16}, with a \emph{d}-spacing of
0.305 nm. On the other hand, for the F\(_{\text{16}}\)CuPc NW sample, we observe
an additional peak corresponding to the (002) plane (labeled as
\(\gamma\)) (Figure \ref{fig:main-2}b right), with its wave vector rotated
by \(\sim 76^{\circ}\) with respect to the \(q_{\mathrm{z}}\) axis and a
\emph{d}-spacing of 1.440 nm, slightly smaller than that in the edge-on
orientation. Together with the (\(\bar{1}\)22) peak located at the
\(q_{\mathrm{z}}\) axis, we point out the emergence of the ``face-on''
orientation (Figure \ref{fig:main-2}b bottom right) in the NWs. The
orientations revealed by GIXD patterns are in consistent with X-ray
diffraction (XRD) experiments (details see Section S2, Figure S10 and
Table S1 in the Supporting Information). The predominant face-on
orientation in the n-type NWs facilitates electron transport
perpendicular to the substrate, through the intermolecular \(\pi-\pi\)
interactions. The substrate-induced orientation transition also
results in the angle-dependent scattering of the Davydov multiplets
that alters the Raman characteristics \cite{Cerdeira_2013_RamanF16}
(details see Supporting Information Figure S11). Figure
\ref{fig:main-2}c presents the scanning electron microscopy (SEM)
micrograph and Raman map giving the intensity ratio of the peak at
1380 cm\(^{\text{-1}}\) to that at 1315 cm\(^{\text{-1}}\), I(1380)/I(1315), on the
F\(_{\text{16}}\)CuPc film deposited on a sheet of SiO\(_{\text{2}}\)-supported graphene
with a break. The NW region, with a higher I(1380)/I(1315) ratio, is
clearly identified.

We examined the wettability of the two F\(_{\text{16}}\)CuPc samples (Figures
\ref{fig:main-2}d-\ref{fig:main-2}g). The apparent static contact
angles (\(\theta_{\mathrm{s}}^{*}\)) of water were independently
determined by (i) a contact angle goniometer (CAG) and (ii) condensation
of  water vapor in an environmental scanning electron
microscope (ESEM) setup. For example, Figure \ref{fig:main-2}d
presents the CAG/ESEM-measured \(\theta_{\mathrm{s}}^{*}\) of a 50 nm-thick
F\(_{\text{16}}\)CuPc/SiO\(_{\text{2}}\) and a 300 nm-thick F\(_{\text{16}}\)CuPc NW samples, showing
103.0\(\pm\)5.4\(^{\circ}\) / 107.0\(\pm\)3.9\(^{\circ}\) and
152.8\(\pm\)3.9\(^{\circ}\) / 149.8\(\pm\)3.4\(^{\circ}\),
respectively. To gain more insights into the superhydrophobic states
of the two surfaces, we conducted a series of static and dynamic
contact angle measurements which allow us to determine
\(\theta_{\mathrm{adv}}^{*}\), \(\theta_{\mathrm{rec}}^{*}\) and
\(\theta_{\mathrm{s}}^{*}\), with respect to the film thickness (Figure
\ref{fig:main-2}e). The apparent contact angle values
\(\theta_{\mathrm{i}}^{*}\), where i=adv, rec, and s, were numerically
fitted by the Wenzel (W) \cite{Wenzel_1936} and CB \cite{Cassie_1944}
models, given by (details see Supplementary Section S3):

\begin{eqnarray}
\label{eq:2}
&\cos \theta^{*}_{\mathrm{i,W}} =& r_{\mathrm{W}} \cos \theta_{\mathrm{i}} \\
&\cos \theta^{*}_{\mathrm{i,CB}} =& r_{\mathrm{CB}} f \cos \theta_{\mathrm{i}} + f - 1
\end{eqnarray}
where \(r_{\mathrm{W}}\) and \(r_{\mathrm{CB}}\) are the roughness ratio
of the wet surface area in the two (W and CB) states,
\(\theta_{\mathrm{i}}\) is the respective contact angles on an ideal
surface and \(f\) is the fraction of solid surface area wet by water in
the CB model. By assuming \(r_{\mathrm{W}} \approx r_{\mathrm{CB}}\)
that linearly increases with film thickness owing to the
intrinsically identical chemical nature of the two surfaces, we find
that the water wettability on F\(_{\text{16}}\)CuPc/SiO\(_{\text{2}}\) and F\(_{\text{16}}\)CuPc NW
samples can be nicely described by the Wenzel and CB models,
respectively (solid lines in Figure \ref{fig:main-2}e), allowing us to
determine \(\theta_{\mathrm{adv}}\)=107\(^{\circ}\),
\(\theta_{\mathrm{rec}}\)=78\(^{\circ}\),
\(\theta_{\mathrm{s}}\)=94\(^{\circ}\), and \emph{f} =0.093 (for 300 nm
NWs) (details see Supporting Information Table S2).  Although some discrepancy
can be seen between the experimental and model data, the simple CB
and Wenzel models still nicely capture the distinct wetting behaviors
on the two F\(_{\text{16}}\)CuPc morphologies. We point out that the
graphene-induced molecular orientation effects have eventually led to
the transition between Wenzel and CB states, which to our knowledge
has never been demonstrated in molecular solids.

We next illustrate the wettability of conductive sessile droplets,
specifically the liquid metals (LMs) at room temperature on the
F\(_{\text{16}}\)CuPc samples. It is noteworthy that the key challenge hindering
the development of LM-based devices is the undesirable adhesion of LM
to the device surface, in particular with the eutectic gallium indium
(EGaIn), in which the formation of gallium oxide layer effectively
reduces the interfacial tension
\cite{Dickey_2008,Doudrick_2014_oxide}. Mercury sessile droplets
sitting on the F\(_{\text{16}}\)CuPc/SiO\(_{\text{2}}\) and F\(_{\text{16}}\)CuPc NW surfaces show
\(\theta_{\mathrm{s}}^{*}\) values of 147.7\(\pm\)2.5\(^{\circ}\) and
157.1\(\pm\)3.8\(^{\circ}\), respectively(Figure \ref{fig:main-2}f).
Regarding the EGaIn sessile droplets, although the
\(\theta_{\mathrm{s}}^{*}\) values of fresh droplets are close on both
surfaces, interestingly, the adhesive property is significantly
different, as revealed in Figure \ref{fig:main-2}g. We suspended an
EGaIn droplet using a microcapillary, followed by repeatedly
touching/removing the droplet from the two surfaces. We found that
the EGaIn droplet can be easily detached from the F\(_{\text{16}}\)CuPc NW
surface (Figure \ref{fig:main-2}g right), compared to that stuck to the
F\(_{\text{16}}\)CuPc/SiO\(_{\text{2}}\) surface (Figure \ref{fig:main-2}g left). The observation was further corroborated
by their apparent sliding angles \(\theta_{\mathrm{sl}}^{*}\),
equivalent to the degree of contact angle hysteresis, of
56.1\(\pm\)13.2\(^{\circ}\) and 14.6\(\pm\)2.6\(^{\circ}\) on
F\(_{\text{16}}\)CuPc/SiO\(_{\text{2}}\) and F\(_{\text{16}}\)CuPc NW samples, respectively (details
see Supplementary Table S3). Accordingly, the excellent CB
characteristics of LM sessile droplets on the semiconducting NWs form
a solid basis for the proposed elastic IFET.

Figure \ref{fig:main-3}a presents the schematic of the proposed elastic
IFET. A sheet of monolayer graphene grown by chemical vapor
deposition (CVD) was transferred onto a 300 nm SiO\(_{\text{2}}\)/Si substrate
functionalized by octadecyltrichrolosilane self-assembled monolayer
(OTS SAM) \cite{Yan_2011}, in order to minimize the substrate-induced
traps \cite{Wang_2011}.  The lateral graphene FET has an on/off current
ratio of \(\sim\)10 and a field-effect mobility (\(\mu_{\mathrm{FE}}\))
of \(\sim\)2000 cm\(^{\text{2}}\)V\(^{\text{-1}}\)s\(^{\text{-1}}\) at room temperature, with the charge
neutrality point (CNP) at the gate voltage of \(\sim\)0 V (Supporting Information Figure S15). Subsequently, a layer of \(\sim\)300 nm-high
F\(_{\text{16}}\)CuPc NWs was deposited on graphene, followed by placing a
cantilever-attached LM droplet on top as the D electrode (details see
Supporting Information Section S1).  The current density
\(J_{\mathrm{DS}}\) from the LM (D) to graphene (S) in the IFET was
then quantified by normalizing the drain current by the contact area
of the sessile droplet with the NW layer, such that \(J_{\mathrm{DS}}
 = {\displaystyle \frac{I_{\mathrm{DS}}}{\pi R^{2}}}\), where \(R\) is
determined from the CAG optical micrograph of the LM droplet (Figure
\ref{fig:main-3}a inset), as a function of drain and gate voltages
(\(V_{\mathrm{D}}\) and \(V_{\mathrm{G}}\)). Note that since the NWs are
only partially wetted, the actual current density passing through
individual NWs is estimated to be \(J_{\mathrm{DS}}/f\). The carrier
density in the F\(_{\text{16}}\)CuPc molecules adjacent to the NW/Gr interface
is modulated by the partially penetrated field effect through
monolayer graphene \cite{Tian_2016,Shih2015Partially}, modulating the
Schottky barrier height at the interface \cite{Yang2012Barristor}. The
face-on orientation of the F\(_{\text{16}}\)CuPc molecules also facilitates
interfacial carrier transport in two ways. The overlap between \(\pi\)
orbitals of F\(_{\text{16}}\)CuPc molecules along the face-on stacking direction
ensures high carrier mobility in the NWs \cite{Bao_1998}. Moreover, the
work function of face-on F\(_{\text{16}}\)CuPc molecules aligns with graphene,
compared with large work function mismatch between edge-on F\(_{\text{16}}\)CuPc
and graphene \cite{Mao_2010_F16_level_orien,Ren_2011_F16Gr_Elevel}. At
zero strain corresponding to the droplet height at \(H_{0}\) determined
analytically (see Equation \ref{eq:H0}, as will be discussed later), the
transfer curve for a representative IFET at \(V_{\mathrm{D}}\)= 1 V
demonstrates an on/off current ratio of 3.9\(\times\)10\(^{\text{4}}\) (Figure
\ref{fig:main-3}b), in line with the state-of-the-art vertical
field-effect transistor (VFET) technology
\cite{Yang2012Barristor,Shih2015Partially,Sun_2017_COF_VFET,Ben_Sasson_2011}. Moreover,
the fact that the present IFET is well-functional at a low
\(V_{\mathrm{D}}\) reflects it potential towards the low-power
electronics design, which remains challenging for most resistive
force-sensing components \cite{Pan_2014,Pang_2012}. A number of IFET
samples were characterized, showing an average on/off current ratio
of 5\(\times\)10\(^{\text{3}}\) at \(V_{\mathrm{D}}\)=1 V (Figure
\ref{fig:main-3}b inset). The transfer curves at different
\(V_{\mathrm{D}}\) are shown in Figure \ref{fig:main-3}c, with the on/off
current ratio gradually decreases with \(V_{\mathrm{D}}\), owing to a
mechanism analogous to the drain-induced barrier lowering (DIBL)
effect in short-channel FETs \cite{Lundstrom_2003}. We note the slight
difference each device may be caused by the air-borne doping of the
IFET, and further efforts will be made to improve the repeatability
of such devices. More discussions about the transport mechanisms,
together with the band diagrams, are stated in Section S4, Figures
S12, S16-S20 in the Supporting Information. Three-terminal operation
of the IFET offers integrated device functionalities sharing with
typical FETs. We demonstrate a circuit composed by an IFET, a green
light-emitting diode (LED), and a bipolar junction transistor (BJT)
amplifier (circuit schematic see Figure \ref{fig:main-3}d left), with
the contact angle monitored by a CAG (Figure \ref{fig:main-3}d
middle). Through continuous forward-reverse scan of \(V_{\mathrm{G}}\),
the total current \(I_{\mathrm{tot}}\) that passes through the circuit,
is modulated by \(\sim\)250 folds and switches the LED on and off,
showing a high current stability and reproducibility (Figure
\ref{fig:main-3}d right). The circuit design maximized the LED response
to compressive stress as well (Supporting Information Movies S1 and S2), as
will be discussed later.

The reversible mechanical response of a CB droplet is the key
principle behind the IFET. To model the elastic response for a CB
droplet, we consider a droplet sandwiched between two flat plates,
with two apparent contact angles \(\theta_{\mathrm{t}}^{*}\) and
\(\theta_{\mathrm{b}}^{*}\), corresponding to the top and bottom
liquid-solid interfaces, respectively. Under the assumptions of (i)
the Bond number Bo\(\ll\)1 and (ii) the contact angles remain
constant independent of \(\sigma\), the cross-sectional boundary of the
droplet can be described as a segment of a perfect circle
\cite{berthier_2012_microdroplet}. Accordingly, the Laplace pressure
\(p\) of the droplet is given by: \(p = \gamma_{\mathrm{L}}
 (R_{1}^{-1} + R_{2}^{-1})\), where \(R_{1}\) and \(R_{2}\) are the
principal radii of the droplet, as schematically shown in Figure
\ref{fig:main-4}a. Upon applying a compressive stress \(\sigma\) between
the plates, the droplet experiences an uniaxial strain \(\varepsilon =
 (H_{0} - H) / H_{0}\), where \(H_{0}\) and \(H\) are the droplet heights
before and after stress, respectively. The compressive stress varies
with height, following \(\sigma = p(H) - p(H_{0})\). Note that here the
liquid phase itself is nearly incompressible, and the ``elasticity'' is
originated from a thermodynamic driving force counteracting the
increase of interfacial tension upon mechanical stress, conceptually
different from the deformation of a bulk material. We formulate the
principal radii as a function of droplet height \(H\) for \(H < H_{0}\),
namely \(R_{1}(H)\) and \(R_{2}(H)\) (detailed derivations see Section S4
and Figure S21 in the Supporting Information). First, the maximum
height \(H_{0}\) corresponding to \(\sigma=0\) is given by:

  \begin{equation}
  \label{eq:H0}
  \begin{aligned}
    H_{0} &= \sqrt[3]{\frac{3 V_{\mathrm{drop}}}{4 \pi}} \sqrt[3]{\frac{1}{ 
   g(\theta_{\mathrm{t}}^{*}) + g(\theta_{\mathrm{b}}^{*}) -1 }}  \left(\cos \theta_{\mathrm{t}}^{*} + \cos \theta_{\mathrm{b}}^{*
}\right) \\
    g(\theta) &= \left(\frac{1 + \cos \theta}{2} \right)^{2} \left(2 - \cos \theta \right)
  \end{aligned}
  \end{equation}
where \(V_{\mathrm{drop}}\) is the the droplet volume following
\(V_{\mathrm{drop}} = w(R_{1}, H, \theta_{\mathrm{t}}^{*},
  \theta_{\mathrm{b}}^{*})\), in which \(w\) is an implicit function of
\(R_{1}\) (see Supporting Information Section S5) that can be solved
numerically for a given \(H\). On the other hand, the second principal
radius is geometrically given by:
\begin{equation}
\label{eq:1}
R_{2} = -\frac{H}{\cos \theta_{\mathrm{t}}^{*} + \cos \theta_{\mathrm{b}}^{*}}
\end{equation}
By using the above equations, the compressive stress \(\sigma\) as a
function of \(\varepsilon\), as well as the effective elastic modulus
\(E = \left({\displaystyle \frac{\mathrm{d} \sigma}{\mathrm{d}
  \varepsilon}}\right)_{H_{0}}\), can be calculated numerically. We
note that \(\theta^{*}_{\mathrm{b}}\) depend on the factor \(f\), rather
than the average radius of the NWs.  To validate our model, a
mercury droplet having an air-stable surface tension
\(\gamma_{\mathrm{L}}\)=0.487 J\(\cdot\)m\(^{\text{-2}}\) is used. Note that
mercury often forms alloys with commonly-used metals
\cite{Kieffer_1959}, so the top contact angle
\(\theta_{\mathrm{t}}^{*}\) may vary depending on the sample
preparation process. For each droplet height, we determined the
experimental \(\sigma\) by extracting the principal radii from the CAG
images, with the height controlled by a stage
micromanipulator. Figure \ref{fig:main-4}b compares the experimental
and calculated elastic stress of a 0.1 \(\mathrm{\mu}\)L droplet as
a function of strain, showing excellent agreement. We notice that
within the strain range considered here (up to 13.5\%), the \(\sigma -
  \varepsilon\) profile is nearly linear, following the Hooke's law. We
determine the effective elastic modulus to be 820 Pa, which is more
than 100 times softer than the widely used Ecoflex rubbers
(\(E\sim{}\)0.1 MPa) \cite{Mosadegh_2014}, and to our knowledge, softer than any solid
materials including the state-of-the-art ultrasoft elastomers
\cite{Miriyev_2017,Jang_2015}. Using our model, we further calculate
\(E\) versus \(V_{\mathrm{drop}}\) (Figure \ref{fig:main-4}b inset,
and Supporting Information Figure S24), predicting an adjustable elastic modulus
by simply controlling the droplet volume. Unsurprisingly, a smaller
droplet tends to be stiffer due to an intrinsically large Laplace
pressure. Another degree of freedom for tunning the elastic modulus
is the bottom contact angle (Supporting Information Figure S24); in other
words, through controlling the wettability of a semiconductor
surface, one may design an IFET with a desirable mechanical
responsiveness.

We next examine the current response of the IFET under mechanical
stress. Figure \ref{fig:main-4}c presents the transfer curves at
different \(\sigma\) values. The current from drain to source
increases with \(\sigma\), with the on/off current ratio unaffected by
the elastic stress. The elastic response at \(V_{\mathrm{G}}\)= 0 V
shows a sensitivity, \(\eta = (I/I_{0} - 1) / \sigma\) , where
\(I_{0}\) and \(I\) are \(I_{\mathrm{DS}}\) before and after stress,
of 36 kPa\(^{-1}\) (Supporting Information Figure S22), together
with the detection limit of down to \textless{}10 Pa, comparable to the
most sensitive resistive pressure sensor
\cite{Mannsfeld_2010,Pang_2012,Pan_2014,Zang_2015}, with a
considerably lower driving voltage that favors low-power-consumption
designs. Following the design concept stated earlier, we point out
that the major mechanism responsible for the current response is an
increased contact area \(\Delta A\) upon stress \(\sigma\), and the
change of penetration length \(\Delta l\) (see Figure \ref{fig:main-1}b)
is negligible within the stress range considered. Indeed, we model
the contact radii at the top and bottom interfaces, \(r_{\mathrm{t}}\)
and \(r_{\mathrm{b}}\), as a function of \(\varepsilon\) (details see
Supporting Information Section S4), which can describe the current
response reasonably well (Supporting Information Figures S22 and
S23). The above analysis is further endorsed by the finite element
method (FEM) simulations (Figure \ref{fig:main-4}d and Supporting Information Table S4).  We further demonstrate the reversibility
of current response by repeatedly applying and releasing a
compressive stress of 204.2\(\pm\)12.3 Pa to our IFET device (Figure
\ref{fig:main-4}e and Supporting Information Movie S3). It is also
noteworthy that the EGaIn sessile droplets function equally well on
the NWs (Supporting Information Figure S19), in spite of their
sticky oxide surface.

Finally, we discuss the transport properties of the NW/Gr
interface. As illustrated earlier, since the carrier density at the
interface is modulated by a partially-penetrated electrostatic field
\cite{Tian_2016,Shih2015Partially}, it has been suggested that the
current density can be described by the thermionic emission model
\cite{Sze2006Mosfets} (details see Section S4 and Figure S25 in the
Supporting Information), which allows to quantify the Schottky
barrier height, \(\Phi_{\mathrm{SB}}\), from the temperature-dependent
measurements. We find that our IFET transfer current substantially
increases with temperature (Figure \ref{fig:main-5}a). The scenario of
temperature-induced increase of contact area or penetration length
is excluded, since the thermal expansion of LM is small (\textless{}1\%)
within the temperature range considered here. Figure \ref{fig:main-5}b
presents the experimentally-obtained current gain as a function of
temperature, \(G(T)=J_{\mathrm{DS}}(T)/J_{\mathrm{DS}}(T=20\ ^{\circ}
  \mathrm{C})\) at different \(V_{\mathrm{G}}\) levels together with the
least-square fitting curves using the thermionic emission
model. Accordingly, at a more negative \(V_{\mathrm{G}}\), the
temperature dependence appears to be stronger, suggesting a higher
\(\Phi_{\mathrm{SB}}\) blocking thermally-induced transport of
electrons. The extracted \(\Phi_{\mathrm{SB}}\) values as a function
of \(V_{\mathrm{G}}\) is shown in Figure \ref{fig:main-5}c, spanning
from 0.46 V at \(V_{\mathrm{G}}\)=-100 V to 0.16 V at
\(V_{\mathrm{G}}\)=100 V.  The range of gate-tunable
\(\Phi_{\mathrm{SB}}\) is comparable to that of the Si/graphene
heterojunction \cite{Yang2012Barristor}, but considerably lower than
our theoretical prediction using the elementary electronic
properties of graphene (Figure \ref{fig:main-5}c and Supporting Information Section S4). A degree of Fermi level pinning due to the
surface-bound traps \cite{Meric_2008} may explain the observation.



In this work, we establish a new concept to reliably incorporate
mechanical durability and responsiveness in a transistor by
engineering the wettability of a semiconductor surface. Using the
superhydrophobic semiconductor NWs in an IFET, we systematically
analyze the origin of its ultrasoft elasticity driven by the
minimization of interfacial tension upon a compressive stress, as well
as the current response that can be modulated by gating through the
atomically-thin 2DEG. We believe that the fundamental understanding
and rational design strategy presented here can be utilized for a wide
range of ultrasensitive and stimuli-responsive nanoelectronics.




\section*{Methods}
\label{sec:orgdf8ba2c}
Detailed descriptions of methods and characterization methods can be seen in the Supporting Information.

\section*{Acknowledgments}
\label{sec:org55f4639}
T.T., N.A., R.S., and C.J.S. are grateful for the startup funding
from ETH Z\"{u}rich. T.T. acknowledges Dr. Tobias Keplinger for
technical support of Raman microscopy. C.S.S. acknowledges Ms. Asel
Maria Aguilar Sanchez from the Institute for Building Materials, ETH
Zurich for her support of ESEM measurements. Y.T.L. and
Y.C.C. acknowledge the Ministry of Science and Technology, Taiwan,
for financial support (project no. MOST 105-2218-E-011-017-MY3).

\section*{Competing Interests}
\label{sec:org91f202c}
The authors declare that they have no competing financial interests.

\section*{Author contributions}
\label{sec:org12c3d57}
T.T. and C.J.S. conceived the concept and designed the
experiments. T.T. and C.S.S. carried out the ESEM
measurements. T.T., N.A. and R.S. fabricated and characterized the
morphology and dynamic wetting properties of the F\(_{\text{16}}\)CuPc
samples. Y.T. L. and Y.C. C. measured the GIXD spectroscopy of
F\(_{\text{16}}\)CuPc samples. T.T. fabricated and tested the interfacial
transistors. T.T. and M.V. designed the experiments using
EGaIn. T.T. and C.J.S. developed the model for the stress-strain
relation of droplets. T.T. and C.J.S. wrote the manuscript. All
authors contributed to the discussion of the results and to the
revision of the manuscript.



% \input{paper.bbl.bak}

\newpage{}

\begin{figure}[htbp]
\centering
% \includegraphics[width=0.95\linewidth]{img/scheme-1.pdf}
\caption{\label{fig:main-1}
Design concept of the proposed elastic interfacial transistor. (a), Schematic showing the device architecture at the micrometer (top) and nanometer (bottom) scales. Superhydrophobic semiconductor nanowires (NWs) are sandwiched between a conductive droplet at the CB state and a layer of 2DEG that allows partial penetration of an electrostatic field. (b), Upon applying an compressive stress \(\sigma\) to the droplet, a vertical displacement \(\Delta H\) advances the contact line by \(\Delta x\), with an apparent advancing contact angle \(\theta_{\mathrm{adv}}^{*}\). When the stress is released, the contact line moves back towards its equilibrium shape, with an apparent receding contact angel of \(\theta_{\mathrm{rec}}^{*}\). At the nanometer scale,  may also result in an increase of penetration length \(\Delta l\) that effectively reduces the transistor channel length.}
\end{figure}


\begin{figure}[htbp]
\centering
% \includegraphics[width=0.9\linewidth]{img/scheme-2.pdf}
\caption{\label{fig:main-2}
Molecular orientation-induced superhydrophobicity in NWs. (a), Top-view (left) and cross-sectional (right) SEM micrographs showing the morphologies of F\(_{\text{16}}\)CuPc/SiO\(_{\text{2}}\) (green) and F\(_{\text{16}}\)CuPc NWs (cyan). Scale bar: 200 nm. (b), GIXD patterns and molecular orientation schematics of F\(_{\text{16}}\)CuPc/SiO\(_{\text{2}}\) (left) and F\(_{\text{16}}\)CuPc NWs (right), showing dominant edge-on (while circles) and face-on (black circles) orientations, respectively. The four representative planes (\(\alpha\) to \(\delta\)) are labeled accordingly. (c), SEM micrograph (top) and Raman mapping image (bottom) of F\(_{\text{16}}\)CuPc molecules deposited onto a graphene sheet with a break, clearly identifying the NW region. Scale bar: 2 \(\mu\)m. (d), Images of water sessile droplets sitting on F\(_{\text{16}}\)CuPc/SiO\(_{\text{2}}\) (top) and F\(_{\text{16}}\)CuPc NWs (bottom) taken from the ESEM and CAG (inset) techniques. (e) CAG-determined water \(\theta_{\mathrm{adv}}^{*}\), \(\theta_{\mathrm{rec}}^{*}\), and \(\theta_{\mathrm{s}}^{*}\) values (dots) with respect to the film thickness (NW height) for the F\(_{\text{16}}\)CuPc/SiO\(_{\text{2}}\) (top) and F\(_{\text{16}}\)CuPc NW (bottom) samples, which are nicely described by the Wenzel and CB models (solid curves), respectively. (f), Mercury sessile droplets sitting on F\(_{\text{16}}\)CuPc/SiO\(_{\text{2}}\) (top) and F\(_{\text{16}}\)CuPc NWs (bottom), showing \(\theta_{\mathrm{s}}^{*}\) values of 147.7\(\pm\)2.5\(^{\circ}\)  and 157.1\(\pm\)3.8\(^{\circ}\), respectively. Scale bar: 500 \(\mu\)m. (g), Touching (left) - removing (right) of a suspended EGaIn droplet from F\(_{\text{16}}\)CuPc/SiO\(_{\text{2}}\) (top) and F\(_{\text{16}}\)CuPc NWs (bottom) surfaces, showing non-stick characteristics of EGaIn to the NWs.}
\end{figure}


\begin{figure}[htbp]
\centering
% \includegraphics[width=0.95\linewidth]{img/scheme-3.pdf}
\caption{\label{fig:main-3}
Transport characteristics of the fabricated IFETs at zero strain. (a), Schematic of the device architecture, together with an optical micrograph taken by CAG. Scale bar: 500 \(\mu\)m. (b), A representative transfer curve at \(V_{\mathrm{D}}\) = 1 V, showing an on/off current ratio of 3.9\(\times{}\)10\(^{\text{4}}\). The inset shows the histogram of the on/off current ratio extracted from \textgreater{}100 IFET devices. (c), Representative transfer curves at different drain voltages. (d), The schematic (left), photographs (middle), and output current \(I_{\mathrm{tot}}\) (right) of an IFET-based circuit that controls the light intensity of a commercial LED with \(V_{\mathrm{G}}\).}
\end{figure}


\begin{figure}[htbp]
\centering
% \includegraphics[width=0.75\linewidth]{img/scheme-4.pdf}
\caption{\label{fig:main-4}
Elastic response of the proposed IFET. (a), Schematics of a droplet sandwiched between two plates before (top) and after (bottom) of a compressive stress, resulting a change of the principal radii, \(R_{1}\) and \(R_{2}\). (b), Compressive stress \(\sigma\) as a function of strain \(\varepsilon\) for a 0.1 \(\mu\)L mercury CB droplet obtained from experiments (dots) and our hydrostatic model (curve). Inset: model-predicted elastic modulus \(E\) as a function of the droplet volume \(V_{\mathrm{drop}}\). (c), Transfer curves under various compressive stress values. Insets: CAG images of the LM droplet under various external stress levels (scale bars: 500 \(\mu\)m.) (d), Comparison of CAG-captured (top) and finite-element simulated (bottom) droplet shape and stress values (Supplementary Table S4) under different degrees of compression. Scale bar: 200 \(\mu\)m. (e), Real-time monitoring of \(I_{\mathrm{DS}}\) by repeatedly applying and releasing a compressive stress.}
\end{figure}


\begin{figure}[htbp]
\centering
% \includegraphics[width=0.95\linewidth]{img/scheme-5.pdf}
\caption{\label{fig:main-5}
Electron transport characteristics at the NW/Gr interface. (a), \(J_{\mathrm{DS}}\) as a function of \(V_{\mathrm{G}}\) at various temperatures, showing the thermoionic emission effect. (b), Experimentally-obtained (dots) and calculated (dashed curves) current gains, \(G(T)\) , as a function of temperature \(T\), suggesting that the thermoionic emission model can describe the interfacial current well. (c), Extracted Schottky barrier height \(\Phi_{\mathrm{SB}}\) (blue dots) as a function \(V_{\mathrm{G}}\). The dashed curve corresponds to our theoretical prediction using the elementary electronic properties of graphene, suggesting a degree of Fermi level pinning at the NW/Gr interface.}
\end{figure}



%%% Local Variables:
%%% mode: latex
%%% TeX-master: "../thesis"
%%% End:
